\documentclass[10pt,a4paper]{amsart}
\usepackage{amsmath,amssymb,amsthm,graphicx}
% For commutative diagrams we use Xy-pic
\usepackage[all]{xy}

% We use 2cell for 2-commutative diagrams.
\xyoption{2cell}
% Theorem environments.
%
\theoremstyle{plain}
\newtheorem{theorem}[subsection]{Theorem}
\newtheorem{proposition}[subsection]{Proposition}
\newtheorem{lemma}[subsection]{Lemma}

\theoremstyle{definition}
\newtheorem{definition}[subsection]{Definition}
\newtheorem{example}[subsection]{Example}
\newtheorem{exercise}[subsection]{Exercise}
\newtheorem{situation}[subsection]{Situation}

\theoremstyle{remark}
\newtheorem{remark}[subsection]{Remark}
\newtheorem{remarks}[subsection]{Remarks}

\numberwithin{equation}{subsection}

% Macros
%
\def\lim{\mathop{\mathrm{lim}}\nolimits}
\def\colim{\mathop{\mathrm{colim}}\nolimits}
\def\Spec{\mathop{\mathrm{Spec}}}
\def\Hom{\mathop{\mathrm{Hom}}\nolimits}
\def\Ext{\mathop{\mathrm{Ext}}\nolimits}
\def\SheafHom{\mathop{\mathcal{H}\!\mathit{om}}\nolimits}
\def\SheafExt{\mathop{\mathcal{E}\!\mathit{xt}}\nolimits}
\def\Sch{\mathit{Sch}}
\def\Mor{\mathop{\mathrm{Mor}}\nolimits}
\def\Ob{\mathop{\mathrm{Ob}}\nolimits}
\def\Sh{\mathop{\mathit{Sh}}\nolimits}
\def\NL{\mathop{N\!L}\nolimits}
\def\CH{\mathop{\mathrm{CH}}\nolimits}
\def\proetale{{pro\text{-}\acute{e}tale}}
\def\etale{{\acute{e}tale}}
\def\QCoh{\mathit{QCoh}}
\def\Ker{\mathop{\mathrm{Ker}}}
\def\Im{\mathop{\mathrm{Im}}}
\def\Coker{\mathop{\mathrm{Coker}}}
\def\Coim{\mathop{\mathrm{Coim}}}

% Boxtimes
%
\DeclareMathSymbol{\boxtimes}{\mathbin}{AMSa}{"02}

%
% Macros for moduli stacks/spaces
%
\def\QCohstack{\mathcal{QC}\!\mathit{oh}}
\def\Cohstack{\mathcal{C}\!\mathit{oh}}
\def\Spacesstack{\mathcal{S}\!\mathit{paces}}
\def\Quotfunctor{\mathrm{Quot}}
\def\Hilbfunctor{\mathrm{Hilb}}
\def\Curvesstack{\mathcal{C}\!\mathit{urves}}
\def\Polarizedstack{\mathcal{P}\!\mathit{olarized}}
\def\Complexesstack{\mathcal{C}\!\mathit{omplexes}}
% \Pic is the operator that assigns to X its picard group, usage \Pic(X)
% \Picardstack_{X/B} denotes the Picard stack of X over B
% \Picardfunctor_{X/B} denotes the Picard functor of X over B
\def\Pic{\mathop{\mathrm{Pic}}\nolimits}
\def\Picardstack{\mathcal{P}\!\mathit{ic}}
\def\Picardfunctor{\mathrm{Pic}}
\def\Deformationcategory{\mathcal{D}\!\mathit{ef}}
\usepackage{cleveref}
\crefname{originalreference}{original reference}{original references}
\Crefname{originalreference}{Original reference}{Original references}
\makeatletter
\newlabel{lemma-category-constructible-abelian}{{09YT}{source}{}{sourcepointer-1}{}}
\newlabel{lemma-category-constructible-abelian@cref}{{[originalreference][0][]09YT}{[1][0][]source}{}{}{}}
\makeatother
\crefname{originalreference}{original reference}{original references}
\Crefname{originalreference}{Original reference}{Original references}
\makeatletter
\newlabel{lemma-category-constructible-modules}{{095N}{source}{}{sourcepointer-2}{}}
\newlabel{lemma-category-constructible-modules@cref}{{[originalreference][0][]095N}{[1][0][]source}{}{}{}}
\makeatother
\crefname{originalreference}{original reference}{original references}
\Crefname{originalreference}{Original reference}{Original references}
\makeatletter
\newlabel{lemma-category-constructible-sets}{{09Y9}{source}{}{sourcepointer-3}{}}
\newlabel{lemma-category-constructible-sets@cref}{{[originalreference][0][]09Y9}{[1][0][]source}{}{}{}}
\makeatother
\crefname{originallemma}{original lemma}{original lemmas}
\Crefname{originallemma}{Original lemma}{Original lemmas}
\makeatletter
\newlabel{sites-lemma-directed-colimits-sections}{{0738}{source}{}{sourcepointer-4}{}}
\newlabel{sites-lemma-directed-colimits-sections@cref}{{[originallemma][0][]0738}{[1][0][]source}{}{}{}}
\makeatother
\title{Constructible sheaves are compact with respect to filtered colimits}
\author{The Stacks Project Authors}
\date{stacks-project:etale-cohomology.tex:local-snapshot-20261011; GNU FDL 1.2 or later}
\begin{document}
\maketitle
\noindent\textit{Editorial note.} All delivered mathematical text is the author's own, copied from the pinned Stacks Project chapter file etale-cohomology.tex; no statement or proof step was rewritten, reordered or removed.\par\medskip
\noindent\textit{Editorial note.} The counted proof is the author's complete argument: the proof is written out in full: the reduction to representable sheaves in the set case with the commutation of sections, and the transport to the abelian and module cases with the n-torsion description.\par\medskip
\noindent\textit{Editorial note.} The label lemma-constructible-is-compact names the counted statement, and the proof environment immediately following it in the source is delivered verbatim as its proof.\par\medskip
\noindent\textit{Editorial note.} The counted statement is carried by the author's own lemma environment in the chapter source.\par\medskip
\noindent\textit{Editorial note.} The statement and proof are copied from the pinned chapter file; the four cited inputs are kept as external pointers with their canonical Stacks tags.\par\medskip
\noindent\textit{Allowed standard prerequisites (editorial).}\begin{itemize}
\item Constructible sheaves of sets, abelian groups and modules as finitely generated colimits. Same chapter, the lemmas describing constructible sheaves of the three kinds as the finitely generated objects built from the representable sheaves and the n-torsion extension sheaves. Supplies the reduction in each of the three cases.
\item Sections commute with filtered colimits. Sites, the lemma computing sections of a filtered colimit of sheaves as the filtered colimit of the sections. Supplies the identity Mor(h\_U, colim G\_i) = colim Mor(h\_U, G\_i) used in the set case.
\end{itemize}\medskip
\ifdefined\hypertarget\hypertarget{sourcepointer-1}{}\fi
\noindent\textit{Source reference.} Original reference 09YT: sources/stacks-etale-cohomology/pinned-original-source.tex, source lines 12471--12471. Prerequisite pointer invoked at proof line(s) 12520; it names a standard fact of the same chapter that the counted argument uses, and every step of the counted proof built on it is retained verbatim in the delivered body.\par\medskip
\ifdefined\hypertarget\hypertarget{sourcepointer-2}{}\fi
\noindent\textit{Source reference.} Original reference 095N: sources/stacks-etale-cohomology/pinned-original-source.tex, source lines 12412--12412. Prerequisite pointer invoked at proof line(s) 12521; it names a standard fact of the same chapter that the counted argument uses, and every step of the counted proof built on it is retained verbatim in the delivered body.\par\medskip
\ifdefined\hypertarget\hypertarget{sourcepointer-3}{}\fi
\noindent\textit{Source reference.} Original reference 09Y9: sources/stacks-etale-cohomology/pinned-original-source.tex, source lines 12343--12343. Prerequisite pointer invoked at proof line(s) 12511; it names a standard fact of the same chapter that the counted argument uses, and every step of the counted proof built on it is retained verbatim in the delivered body.\par\medskip
\ifdefined\hypertarget\hypertarget{sourcepointer-4}{}\fi
\noindent\textit{Source reference.} Original lemma 0738: sources/stacks-chow/tags.txt, source lines 9052--9052. Prerequisite pointer invoked at proof line(s) 12515; it names a standard fact of a sibling chapter that the counted argument uses, and every step of the counted proof built on it is retained verbatim in the delivered body.\par\medskip
\begin{lemma}
\label{lemma-constructible-is-compact}
Let $X$ be a quasi-compact and quasi-separated scheme. Let $\Lambda$ be a
Noetherian ring. Let $\mathcal{F}$ be a constructible sheaf of sets, abelian
groups, or $\Lambda$-modules on $X_\etale$. Let
$\mathcal{G} = \colim \mathcal{G}_i$ be a filtered colimit of sheaves of
sets, abelian groups, or $\Lambda$-modules. Then
$$
\Mor(\mathcal{F}, \mathcal{G}) = \colim \Mor(\mathcal{F}, \mathcal{G}_i)
$$
in the category of sheaves of sets, abelian groups, or $\Lambda$-modules on
$X_\etale$.
\end{lemma}

\begin{proof}
The case of sheaves of sets. By Lemma \ref{lemma-category-constructible-sets}
it suffices to prove the lemma for $h_U$ where $U$ is a quasi-compact
and quasi-separated object of $X_\etale$. Recall that
$\Mor(h_U, \mathcal{G}) = \mathcal{G}(U)$. Hence the result
follows from Sites, Lemma \ref{sites-lemma-directed-colimits-sections}.

\medskip\noindent
In the case of abelian sheaves or sheaves of modules, the result follows
in the same way using
Lemmas \ref{lemma-category-constructible-abelian} and
\ref{lemma-category-constructible-modules}.
For the case of abelian sheaves, we add that
$\Mor(j_{U!}\underline{\mathbf{Z}/n\mathbf{Z}}, \mathcal{G})$
is equal to the $n$-torsion elements of $\mathcal{G}(U)$.
\end{proof}

\end{document}
